\chapter{The Mean Value Theorem and What It Means}

There is a moment in every calculus course where a single, almost obvious-sounding
fact suddenly turns out to be doing all the heavy lifting behind the scenes. You
learn that a function with positive derivative is increasing. You learn that two
antiderivatives of the same function differ by a constant. You learn that you can
bound the error of a linear approximation. Each of these feels like its own little
truth, earned separately. They are not. They are all the same theorem, wearing
different coats.

That theorem is the \emph{Mean Value Theorem} (MVT). Once you \emph{see} why it is
true, you will start noticing it quietly powering a large fraction of everything
that follows. The MVT is the bridge that carries
information from the derivative back to the function. Up to now we have spent our
energy going one direction --- from $f$ to $f'$. The MVT is the first serious tool
for going the other way: knowing something about $f'$ on an interval, what can we
conclude about $f$ itself? Almost everything, it turns out.

The build happens in three stages: a picture so simple it feels like a triviality
(Rolle's theorem), the small tilt that turns it into the full MVT, and then a long,
satisfying tour of the consequences, where the ``rules'' you half-trusted turn out
to be theorems with real proofs.

\section{The picture you already believe}

Take a smooth curve and draw the straight line connecting its two endpoints. That
line --- the \emph{secant} --- has some slope, the average rate of change of the
function across the interval. Now slide a ruler parallel to that secant, keeping it
parallel, and push it up against the curve until it just kisses it. The point where
it touches is a point whose \emph{tangent} has exactly the slope of the secant.

\begin{intuitionbox}
The Mean Value Theorem says: somewhere on a smooth arc, the instantaneous slope
equals the average slope. If you drove $120$ miles in $2$ hours, then at some
instant your speedometer read exactly $60$ mph --- not more, not less, exactly the
average. You cannot average $60$ without \emph{being} at $60$ at some moment. That
is the entire content of the theorem, and it is hard to believe it could be the
foundation of so much.
\end{intuitionbox}

\begin{figure}[H]
\centering
\begin{tikzpicture}
\begin{axis}[
  width=11cm, height=7cm,
  axis lines=middle,
  xlabel={$x$}, ylabel={$y$},
  xtick={1,4}, xticklabels={$a$,$b$},
  ytick=\empty,
  xmin=0, xmax=5.2, ymin=-0.4, ymax=3.6,
  samples=120, domain=0.4:4.8,
  clip=false,
]
% the curve f(x)=0.25*(x-2.3)^2 + 1 + 0.35*x
\addplot[curveblue, very thick, domain=1:4] {0.25*(x-2.3)^2 + 1 + 0.35*x};
% endpoints: f(1)=1.7725, f(4)=3.1225
\addplot[only marks, mark=*, mark size=1.6pt, black]
  coordinates {(1,1.7725) (4,3.1225)};
% secant line through endpoints: slope = (3.1225-1.7725)/3 = 0.45
\addplot[tangentred, thick, domain=0.4:4.8] {1.7725 + 0.45*(x-1)};
% tangent point c where f'(x)=0.5(x-2.3)+0.35=0.45 -> x-2.3=0.2 -> x=2.5, f(2.5)=1.885
\addplot[tangentred, thick, dashed, domain=1.3:3.7]
  {1.885 + 0.45*(x-2.5)};
\addplot[only marks, mark=*, mark size=1.8pt, tangentred]
  coordinates {(2.5, 1.885)};
\node[tangentred] at (axis cs:2.5,1.45) {$c$};
\draw[dashed, gray] (axis cs:2.5,1.885) -- (axis cs:2.5,0);
\node[curveblue] at (axis cs:4.3,3.2) {$y=f(x)$};
\node[tangentred] at (axis cs:0.9,2.55) {secant};
\end{axis}
\end{tikzpicture}
\caption{The Mean Value Theorem. The red secant joins the endpoints
$(a,f(a))$ and $(b,f(b))$. Somewhere in between, at $x=c$, the tangent (dashed) is
parallel to it: $f'(c)$ equals the average slope.}
\label{ch15:fig-mvt}
\end{figure}

Notice what the picture quietly assumes. The curve has to be unbroken --- if it
jumped, the ruler could miss it entirely. And it has to be smooth, with a
well-defined tangent everywhere inside, or there could be a sharp corner where the
slope you want is never achieved. Hold onto those two requirements; they will
become the hypotheses of the theorem, and we will see exactly why each one is
needed.

\section{Rolle's theorem: the flat special case}

The cleanest way into the MVT is to first prove a special case so plain it almost
proves itself, and then \emph{tilt} the whole picture to get the general statement.
The special case is the one where the two endpoints sit at the same height.

\begin{intuitionbox}
If a smooth curve starts and ends at the same height, then somewhere in between it
has a horizontal tangent. Throw a ball; it leaves your hand and returns to the same
height; at the top of its arc, for one instant, it is moving neither up nor down.
That instant is Rolle's theorem.
\end{intuitionbox}

\begin{theorem}[Rolle, 1691]\label{ch15:thm-rolle}
Let $f$ be continuous on the closed interval $[a,b]$ and differentiable on the open
interval $(a,b)$, and suppose $f(a)=f(b)$. Then there exists at least one point
$c\in(a,b)$ with
\[
  f'(c)=0.
\]
\end{theorem}

\begin{figure}[H]
\centering
\begin{tikzpicture}
\begin{axis}[
  width=10cm, height=5.5cm,
  axis lines=middle,
  xlabel={$x$}, ylabel={$y$},
  xtick={1,4}, xticklabels={$a$,$b$},
  ytick=\empty,
  xmin=0, xmax=5, ymin=-0.3, ymax=2.6,
  samples=120,
  clip=false,
]
% f(x)=1.5+0.55*sin((x-1)/3*pi): equal ends f(1)=f(4)=1.5, max at x=2.5, f=2.05
\addplot[curveblue, very thick, domain=1:4] {1.5 + 0.55*sin((x-1)/3*180)};
\addplot[only marks, mark=*, mark size=1.6pt, black] coordinates {(1,1.5) (4,1.5)};
% horizontal tangent at the max x=2.5, y=2.05
\addplot[tangentred, thick, domain=1.7:3.3] {2.05};
\node[tangentred] at (axis cs:2.5,2.32) {$f'(c)=0$};
\draw[dashed,gray] (axis cs:2.5,2.05)--(axis cs:2.5,0);
\node[gray] at (axis cs:2.5,-0.18) {$c$};
\addplot[only marks, mark=*, mark size=1.6pt, tangentred] coordinates {(2.5,2.05)};
\end{axis}
\end{tikzpicture}
\caption{Rolle's theorem: equal endpoint heights force a horizontal tangent
somewhere inside.}
\label{ch15:fig-rolle}
\end{figure}

\subsection{Why it has to be true}

Here is the idea before the proof, because the idea is the whole point. A
continuous function on a closed interval has to attain a highest value and a lowest
value somewhere --- it cannot run off to infinity, and it cannot sneak up to a
maximum without reaching it, because there are no gaps. (This is the Extreme Value
Theorem, the companion of the Intermediate Value Theorem from the continuity
chapter.) Now, the endpoints are tied together at the same height. So at least one
of the highest or lowest point must occur \emph{strictly inside} the interval ---
unless the function is perfectly flat, in which case every interior point works
anyway. And at an interior peak or valley, a differentiable function must have
derivative zero. Why? Because at a peak the function isn't rising on either side,
so its slope can be neither positive nor negative. The only number that is neither
positive nor negative is zero.

\begin{whybox}
The hidden engine of Rolle's theorem is a tiny fact worth isolating, sometimes
called \emph{Fermat's theorem on interior extrema}: \textit{if $f$ has a local
maximum or minimum at an interior point $c$ and $f'(c)$ exists, then $f'(c)=0$.}

Watch the mechanism. Suppose $c$ is a local max. For $h>0$ small, $f(c+h)\le f(c)$,
so the right-hand difference quotient $\frac{f(c+h)-f(c)}{h}\le 0$; its limit gives
$f'(c)\le 0$. For $h<0$ small, $f(c+h)\le f(c)$ still, but now we divide by a
\emph{negative} $h$, flipping the inequality: $\frac{f(c+h)-f(c)}{h}\ge 0$, so
$f'(c)\ge 0$. The derivative exists, so both one-sided limits equal it. A number
that is both $\le 0$ and $\ge 0$ is $0$. The two sides of the peak squeeze the
slope to zero from opposite directions.
\end{whybox}

\begin{proof}[Proof of Rolle's theorem]
Since $f$ is continuous on the closed, bounded interval $[a,b]$, the Extreme Value
Theorem guarantees that $f$ attains a maximum value $M$ and a minimum value $m$ on
$[a,b]$.

\emph{Case 1: $M=m$.} Then $f$ is constant, $f'(x)=0$ for every $x\in(a,b)$, and any
interior point serves as $c$.

\emph{Case 2: $M>m$.} At least one of these two values differs from the common
endpoint height $f(a)=f(b)$. Say it is the maximum $M$ (the minimum case is
identical). Then the point where $f$ attains $M$ cannot be an endpoint, so it is
some interior point $c\in(a,b)$. At $c$, $f$ has a local maximum, $f$ is
differentiable, and by Fermat's interior-extremum fact $f'(c)=0$.
\end{proof}

\begin{trapbox}
\textbf{Every hypothesis earns its keep.} Drop any one and Rolle fails:
\begin{itemize}
\item \emph{Continuity on $[a,b]$.} Let $f(x)=x$ on $[0,1)$ and redefine $f(1)=0$.
Endpoints match ($f(0)=f(1)=0$), but the derivative is $1$ everywhere it exists ---
no horizontal tangent. The jump at the right end let the function ``escape''.
\item \emph{Differentiability on $(a,b)$.} Let $f(x)=\abs{x}$ on $[-1,1]$. Endpoints
match, $f$ is continuous, but the only place the slope could be zero is the corner
at $0$, where $f'$ does not exist. The peak is there, but it is a spike, not a
smooth top.
\item \emph{Equal endpoints.} Without $f(a)=f(b)$ there is no reason for a flat
spot --- that is precisely the gap the full MVT fills next.
\end{itemize}
\end{trapbox}

\section{Tilting Rolle into the Mean Value Theorem}

Rolle's theorem is the MVT for a level interval. To get the general statement, we
play a trick that is one of the most useful moves in all of analysis: we
\emph{subtract off the secant line} so that the tilted problem becomes a level one.

\begin{intuitionbox}
Imagine the MVT picture drawn on a tilted sheet of paper. Now rotate your head so
the secant looks horizontal. In that tilted view, the curve starts and ends at the
same height --- it is exactly a Rolle picture! Rolle hands us a horizontal tangent
\emph{in the tilted frame}, which is a tangent parallel to the secant \emph{in the
original frame}. We do the ``rotation'' algebraically by subtracting the secant's
formula from $f$.
\end{intuitionbox}

\begin{theorem}[Mean Value Theorem, Lagrange form]\label{ch15:thm-mvt}
Let $f$ be continuous on $[a,b]$ and differentiable on $(a,b)$. Then there exists at
least one point $c\in(a,b)$ such that
\[
  f'(c)=\frac{f(b)-f(a)}{b-a}.
\]
Equivalently, $f(b)-f(a)=f'(c)\,(b-a)$ for some $c$ strictly between $a$ and $b$.
\end{theorem}

\begin{proof}
The secant line through the endpoints is
\[
  L(x)=f(a)+\frac{f(b)-f(a)}{b-a}\,(x-a),
\]
a straight line, so $L'(x)=\dfrac{f(b)-f(a)}{b-a}$ everywhere. Define the
\emph{gap function}
\[
  g(x)=f(x)-L(x),
\]
the vertical distance from the curve to its own secant. Then:
\begin{itemize}
\item $g$ is continuous on $[a,b]$ and differentiable on $(a,b)$, being a difference
of two such functions.
\item $g(a)=f(a)-L(a)=f(a)-f(a)=0$, and similarly $g(b)=f(b)-L(b)=0$. The curve
meets its secant at the endpoints, so the gap is zero there.
\end{itemize}
Thus $g$ satisfies the hypotheses of Rolle's theorem with $g(a)=g(b)=0$. Rolle
gives a point $c\in(a,b)$ with $g'(c)=0$. But
\[
  g'(x)=f'(x)-L'(x)=f'(x)-\frac{f(b)-f(a)}{b-a},
\]
so $g'(c)=0$ means exactly $f'(c)=\dfrac{f(b)-f(a)}{b-a}$.
\end{proof}

\begin{happeningbox}
What is $g(x)=f(x)-L(x)$ really? It is the signed vertical gap between the curve and
the straight line joining its ends. The proof says: the place where this gap is
widest (or narrowest) --- where $g'=0$ --- is the place where the curve is, for an
instant, running \emph{parallel} to the secant. That is geometrically inevitable.
As you move along the curve away from the secant and then back to it, there has to
be a turning point in the gap, and at that turning point the curve and the secant
are momentarily heading the same direction. The MVT is not really about
derivatives at all; it is about the fact that a gap that starts and ends at zero
must, somewhere, stop growing.
\end{happeningbox}

\begin{figure}[H]
\centering
\begin{tikzpicture}
\begin{axis}[
  width=11cm, height=6.5cm,
  axis lines=middle,
  xlabel={$x$}, ylabel={$y$},
  xtick={1,4}, xticklabels={$a$,$b$},
  ytick=\empty,
  xmin=0, xmax=5.2, ymin=-0.4, ymax=3.4,
  samples=140, domain=1:4,
  clip=false,
]
\addplot[curveblue, very thick] {0.25*(x-2.3)^2 + 1 + 0.35*x};
\addplot[tangentred, thick] {1.7725 + 0.45*(x-1)};
% the gap g(x) drawn as vertical segments
\addplot[gray, dashed, samples at={1.4,1.8,2.5,3.0,3.6},
  ycomb, mark=none, thick] {0.25*(x-2.3)^2 + 1 + 0.35*x};
\addplot[gray, dashed, samples at={1.4,1.8,2.5,3.0,3.6},
  ycomb, mark=none, thick] {1.7725 + 0.45*(x-1)};
\node[curveblue] at (axis cs:4.45,3.0) {$f$};
\node[tangentred] at (axis cs:4.55,2.45) {secant $L$};
\node[gray] at (axis cs:2.9,2.55) {gap $g=f-L$};
\end{axis}
\end{tikzpicture}
\caption{The proof in one image: $g=f-L$ is the gap between the curve and its
secant. It vanishes at both ends, so by Rolle it has a horizontal tangent
somewhere --- and there $f$ runs parallel to $L$.}
\label{ch15:fig-gap}
\end{figure}

\begin{historybox}
The names attached here are a little unfair to history. \textbf{Michel Rolle}
(1652--1719) was a French algebraist who, ironically, was for a time a fierce critic
of the new calculus, calling it a collection of ``ingenious fallacies.'' His
theorem appeared in a 1691 book on equations, stated for polynomials, with no
mention of derivatives in the modern sense. The general Mean Value Theorem belongs
to the era of \textbf{Joseph-Louis Lagrange} (1736--1813), who used it as a
foundation stone in his attempt to build calculus on algebra without limits.

The proof you just read --- subtract the secant, apply Rolle --- and the
\emph{rigorous} version of the whole story is due to \textbf{Augustin-Louis Cauchy}
(1789--1857). The early users of calculus, from Newton and Leibniz onward,
\emph{believed} the MVT intuitively (of course an average speed is achieved at some
instant!), and they used it freely. What they lacked was the Extreme Value Theorem
and a clear notion of continuity to justify the one delicate step: that a maximum is
actually attained. Cauchy and later Weierstrass supplied exactly that missing
floor. The MVT is a perfect case study in how calculus was \emph{used} correctly for
two centuries before anyone could fully \emph{prove} it.
\end{historybox}

\subsection{Reading the theorem carefully}

The MVT is an \emph{existence} theorem. It promises that such a point $c$ exists; it
does not tell you where, and there may be several. For most purposes this is
exactly enough, because we will use it to say ``$f'$ at some unknown point equals
this average,'' and then bound or sign that unknown derivative.

\begin{focusbox}
\textbf{The two faces of the MVT.} Keep both forms in mind; you will reach for
different ones in different situations.
\begin{itemize}
\item \emph{Slope form:} $\displaystyle f'(c)=\frac{f(b)-f(a)}{b-a}$. Read it as
``the average rate is hit exactly.''
\item \emph{Increment form:} $f(b)-f(a)=f'(c)(b-a)$. Read it as ``the total change
equals some interior rate times the elapsed input.'' This is the workhorse form:
it converts a statement about $f'$ into a statement about $f(b)-f(a)$.
\end{itemize}
\end{focusbox}

\begin{example}\label{ch15:ex-find-c}
Let $f(x)=x^3$ on $[0,2]$. The average slope is
$\frac{f(2)-f(0)}{2-0}=\frac{8-0}{2}=4$. The MVT promises a $c\in(0,2)$ with
$f'(c)=3c^2=4$, i.e. $c^2=\frac43$, $c=\frac{2}{\sqrt3}\approx 1.155$. Indeed
$1.155$ lies inside $(0,2)$, as advertised. Notice we \emph{found} $c$ here only
because the function is simple; in real applications we almost never solve for $c$
--- we just use the fact that some such $c$ exists.
\end{example}

\begin{labbox}
Take $f(x)=\sqrt{x}$ on $[1,4]$. Compute the average slope
$\frac{\sqrt4-\sqrt1}{4-1}=\frac{2-1}{3}=\frac13$. Now set $f'(c)=\frac{1}{2\sqrt
c}=\frac13$ and solve: $2\sqrt c=3$, $\sqrt c=\frac32$, $c=\frac94=2.25$. Check that
$2.25$ sits in $(1,4)$. You have just located the exact instant where this curve's
slope equals its average climb. Try the same with $f(x)=1/x$ on $[1,2]$ and confirm
$c=\sqrt2$.
\end{labbox}

\section{Why the MVT is the workhorse}

Now the payoff. Everything in this section follows from the single increment form
$f(b)-f(a)=f'(c)(b-a)$. The strategy is always the same: \emph{whatever you know
about $f'$ on the interval, you know about that mystery factor $f'(c)$, and the
equation passes the knowledge on to $f(b)-f(a)$.} Watch how many ``obvious'' facts
turn out to be this one theorem in disguise.

\subsection{Zero derivative means constant}

You believe that if a function never changes (its rate is always $0$) then it stays
put. But why, \emph{exactly}? ``The derivative is zero'' is a statement about
slopes; ``the function is constant'' is a statement about values. The MVT is the
only rigorous bridge between them.

\begin{theorem}[Vanishing derivative]\label{ch15:thm-const}
If $f$ is continuous on $[a,b]$, differentiable on $(a,b)$, and $f'(x)=0$ for all
$x\in(a,b)$, then $f$ is constant on $[a,b]$.
\end{theorem}

\begin{proof}
Pick any point $x\in(a,b]$. Apply the MVT to $f$ on the subinterval $[a,x]$: there
is a $c$ between $a$ and $x$ with
\[
  f(x)-f(a)=f'(c)\,(x-a).
\]
But $f'(c)=0$ by hypothesis, so $f(x)-f(a)=0$, i.e. $f(x)=f(a)$. Since $x$ was
arbitrary, $f$ equals the constant $f(a)$ everywhere.
\end{proof}

\begin{trapbox}
This is exactly the place where the \emph{interval matters}. On a disconnected
domain the conclusion is false. The function
\[
  f(x)=\begin{cases}1,& x>0,\\ -1,& x<0,\end{cases}
\]
has $f'(x)=0$ everywhere it is defined, yet it is not constant. There is no
contradiction: the MVT needs a single interval connecting any two points, and there
is no such interval joining a negative $x$ to a positive one inside the domain. The
derivative tells you nothing about how $f$ behaves \emph{across a gap}.
\end{trapbox}

\subsection{Two antiderivatives differ by a constant}

This is the fact that makes the ``$+\,C$'' in every indefinite integral legitimate.
If you have found \emph{one} antiderivative of a function, you have essentially
found them all.

\begin{corollary}[Antiderivatives up to a constant]\label{ch15:cor-plusC}
If $F'(x)=G'(x)$ for all $x$ in an interval $I$, then $F(x)=G(x)+C$ on $I$ for some
constant $C$.
\end{corollary}

\begin{proof}
Let $h=F-G$. Then $h'(x)=F'(x)-G'(x)=0$ for all $x\in I$. By the previous theorem,
$h$ is constant on $I$: $h(x)=C$. Hence $F(x)=G(x)+C$.
\end{proof}

\begin{aestheticbox}
This little corollary is quietly enormous. It is the reason the Fundamental Theorem
of Calculus can speak of \emph{the} antiderivative up to a constant; the reason
$\int f$ is well posed; the reason solving a differential equation $y'=f$ gives a
one-parameter family rather than chaos. The entire bookkeeping of integration ---
that mysterious $+C$ you were told to ``never forget'' --- is the MVT speaking. You
were not memorizing a rule. You were memorizing a theorem.
\end{aestheticbox}

\subsection{The increasing/decreasing test}

Here is the workhorse role most people meet first without realizing where it comes
from. We \emph{define} ``increasing'' by values: $f$ is increasing on $I$ if
$x_1<x_2$ forces $f(x_1)<f(x_2)$. The sign of the derivative is, again, a statement
about slopes. The MVT connects the two.

\begin{theorem}[Monotonicity test]\label{ch15:thm-mono}
Let $f$ be continuous on an interval $I$ and differentiable in its interior.
\begin{itemize}
\item If $f'(x)>0$ throughout the interior, then $f$ is strictly increasing on $I$.
\item If $f'(x)<0$ throughout, then $f$ is strictly decreasing on $I$.
\item If $f'(x)\ge 0$ throughout, then $f$ is nondecreasing; if $f'(x)\le 0$,
nonincreasing.
\end{itemize}
\end{theorem}

\begin{proof}
Take any two points $x_1<x_2$ in $I$. By the MVT on $[x_1,x_2]$ there is a $c$
between them with
\[
  f(x_2)-f(x_1)=f'(c)\,(x_2-x_1).
\]
The factor $(x_2-x_1)$ is positive. So the sign of $f(x_2)-f(x_1)$ is exactly the
sign of $f'(c)$. If $f'>0$ everywhere, then $f'(c)>0$, so $f(x_2)-f(x_1)>0$, i.e.
$f(x_2)>f(x_1)$: increasing. The other three cases are read off the same equation by
changing the sign or allowing equality.
\end{proof}

\begin{whybox}
\textbf{Why does the \emph{sign} of the derivative control the \emph{direction} of
the function?} Because of that single MVT equation. Without it, ``$f'>0$'' and
``$f$ increasing'' are two unrelated sentences --- one about tangent lines, one
about a table of values. The MVT says the total change $f(x_2)-f(x_1)$ is built out
of an actual interior slope times a positive width. You cannot accumulate a net
\emph{decrease} out of slopes that are all positive. The local information (slope at
points) integrates up into global information (rise across an interval), and the MVT
is the gear that meshes the two.
\end{whybox}

\begin{trapbox}
\textbf{A subtle but important converse warning.} ``Strictly increasing'' does
\emph{not} require $f'>0$ at every point. Consider $f(x)=x^3$. It is strictly
increasing on all of $\R$, yet $f'(0)=0$. The theorem's arrow points one way:
$f'>0\Rightarrow$ increasing, but increasing only gives $f'\ge 0$. A single flat
instant is allowed, as long as the function does not actually pause over a whole
subinterval. Many students ``prove'' a function is increasing by checking $f'\ge 0$
and then wrongly assert it must be \emph{strictly} increasing; that needs the extra
observation that $f'$ is not zero on any subinterval.
\end{trapbox}

\begin{figure}[H]
\centering
\begin{tikzpicture}
\begin{axis}[
  width=12cm, height=5cm,
  axis lines=middle,
  xlabel={$x$}, ylabel={$y$},
  ytick=\empty, xtick=\empty,
  xmin=-2.6, xmax=2.6, ymin=-1.2, ymax=3.2,
  samples=160, domain=-2.3:2.3,
  clip=false,
]
\addplot[curveblue, very thick] {0.35*x^3 - 1.2*x + 1};
% mark increasing/decreasing regions: f'=1.05x^2-1.2=0 -> x=+-1.069
\draw[dashed,gray] (axis cs:-1.069,-1.2)--(axis cs:-1.069,3.2);
\draw[dashed,gray] (axis cs:1.069,-1.2)--(axis cs:1.069,3.2);
\node[targetgreen] at (axis cs:-1.8,2.6) {$f'>0$};
\node[tangentred] at (axis cs:0,2.6) {$f'<0$};
\node[targetgreen] at (axis cs:1.8,2.6) {$f'>0$};
\node[gray] at (axis cs:-1.069,-1.0) {$c_1$};
\node[gray] at (axis cs:1.069,-1.0) {$c_2$};
\end{axis}
\end{tikzpicture}
\caption{The sign of $f'$ reads off the direction of $f$. Where $f'>0$ the curve
climbs; where $f'<0$ it falls; the turning points $c_1,c_2$ are where $f'=0$. The
monotonicity test is just the MVT applied on each piece.}
\label{ch15:fig-mono}
\end{figure}

\subsection{Error bounds: the MVT as a ruler}

Perhaps the most \emph{useful} face of the MVT in applied work is its power to
\emph{bound} how much a function can change. If you know the speedometer never
exceeds $70$, you know that in two hours you went at most $140$ miles. That is the
MVT used as an inequality.

\begin{corollary}[Change is bounded by the steepest slope]\label{ch15:cor-bound}
Suppose $f$ is differentiable on an interval and $\abs{f'(x)}\le M$ for all $x$
there. Then for any two points $a,b$ in the interval,
\[
  \abs{f(b)-f(a)}\le M\,\abs{b-a}.
\]
\end{corollary}

\begin{proof}
By the MVT, $f(b)-f(a)=f'(c)(b-a)$ for some interior $c$. Take absolute values:
$\abs{f(b)-f(a)}=\abs{f'(c)}\,\abs{b-a}\le M\abs{b-a}$, since $\abs{f'(c)}\le M$.
\end{proof}

A function whose derivative is bounded by $M$ in this way is called
\emph{Lipschitz} with constant $M$: it cannot change faster than $M$ units of output
per unit of input, ever. This single inequality is the backbone of error analysis,
numerical methods, and the existence-and-uniqueness theory for differential
equations.

\begin{example}[Bounding a value of sine]\label{ch15:ex-sin-bound}
How big can $\sin(0.1)$ be, without a calculator? Apply the bound to $f=\sin$ on
$[0,0.1]$. We have $f'=\cos$, and $\abs{\cos x}\le 1$ everywhere, so $M=1$. Then
\[
  \abs{\sin(0.1)-\sin(0)}\le 1\cdot\abs{0.1-0}=0.1.
\]
Thus $\abs{\sin(0.1)}\le 0.1$. We just proved the famous inequality
$\abs{\sin x}\le \abs{x}$ for all $x$ --- which is normally drawn from a picture ---
as a one-line consequence of the MVT, because $\abs{(\sin)'}=\abs{\cos}\le 1$.
\end{example}

\begin{example}[A sharper estimate near a known value]\label{ch15:ex-sqrt-bound}
Estimate $\sqrt{4.1}$ and bound the error of the estimate $2.025$. With
$f(x)=\sqrt x$ on $[4,4.1]$, the MVT gives
\[
  \sqrt{4.1}-\sqrt 4=f'(c)(0.1)=\frac{1}{2\sqrt c}(0.1)
  \quad\text{for some }c\in(4,4.1).
\]
Now $c$ is somewhere in $(4,4.1)$, so $\frac{1}{2\sqrt c}$ lies between
$\frac{1}{2\sqrt{4.1}}$ and $\frac{1}{2\sqrt 4}=0.25$. Hence
\[
  \sqrt{4.1}=2+\frac{0.1}{2\sqrt c}, \qquad
  2+0.1(0.2469)\le \sqrt{4.1}\le 2+0.1(0.25),
\]
so $2.02469\le \sqrt{4.1}\le 2.025$. The MVT didn't just give a guess; it
\emph{trapped} the true value between two computable numbers. This is the engine
behind rigorous numerical bounds.
\end{example}

\begin{applicationbox}
\textbf{Where this bites in practice.} In numerical analysis, every statement of the
form ``this approximation is accurate to within $\epsilon$'' is ultimately an MVT
estimate: the error is some derivative (or higher derivative, via Taylor's theorem,
which is the MVT applied repeatedly) evaluated at an unknown interior point, and you
bound it by the worst-case size of that derivative. In machine learning, a loss
function with a derivative bounded by $M$ (an $M$-Lipschitz gradient) is exactly the
condition that gradient descent with a small enough step size cannot overshoot ---
the same Corollary~\ref{ch15:cor-bound} controlling how far a step can move you. In
control and physics, a bounded force means a bounded change in velocity over a
bounded time, again the very same inequality.
\end{applicationbox}

\subsection{Proving inequalities without a single graph}

The MVT turns many ``stare at the graph and believe it'' inequalities into clean
proofs. The recipe: build a function that is the difference of the two sides, show
its derivative has a fixed sign, and let monotonicity finish the job.

\begin{example}[$e^x \ge 1+x$ for all $x$]\label{ch15:ex-exp-ineq}
Let $g(x)=e^x-(1+x)$. Then $g(0)=0$ and $g'(x)=e^x-1$. Now $g'(x)>0$ for $x>0$ and
$g'(x)<0$ for $x<0$. By the monotonicity test, $g$ decreases up to $0$ and
increases after it, so $x=0$ is a global minimum of $g$, with value $g(0)=0$. Hence
$g(x)\ge 0$ everywhere, i.e. $e^x\ge 1+x$. The tangent line $1+x$ lies below the
exponential (Figure~\ref{ch15:fig-exp-ineq}), proved purely from the sign of a
derivative.
\end{example}

\begin{figure}[H]
\centering
\begin{tikzpicture}
\begin{axis}[
  width=10cm, height=6cm,
  axis lines=middle,
  xlabel={$x$}, ylabel={$y$},
  xmin=-2.2, xmax=1.6, ymin=-1.3, ymax=4.6,
  samples=120, domain=-2.1:1.5,
  legend pos=north west, legend cell align=left,
  clip=false,
]
\addplot[curveblue, very thick] {exp(x)};
\addlegendentry{$e^x$}
\addplot[tangentred, thick] {1+x};
\addlegendentry{$1+x$}
\addplot[only marks, mark=*, mark size=1.4pt, black] coordinates {(0,1)};
\node at (axis cs:0.55,0.7) {touch at $x=0$};
\end{axis}
\end{tikzpicture}
\caption{The tangent line $y=1+x$ at $x=0$ lies entirely below $y=e^x$, touching
only at the origin. The single fact $g'(x)=e^x-1$ changes sign exactly at $0$
forces this: $g=e^x-(1+x)$ falls to its minimum at $0$ and climbs away on both
sides.}
\label{ch15:fig-exp-ineq}
\end{figure}

\begin{example}[$\ln(1+x)<x$ for $x>0$]\label{ch15:ex-log-ineq}
Let $g(x)=x-\ln(1+x)$ on $[0,\infty)$. Then $g(0)=0$ and
$g'(x)=1-\frac{1}{1+x}=\frac{x}{1+x}>0$ for $x>0$. So $g$ is strictly increasing on
$[0,\infty)$, giving $g(x)>g(0)=0$, that is $\ln(1+x)<x$. Same recipe, opposite
function.
\end{example}

\begin{walkbox}
Here is a question to chew on next time you are out walking. The monotonicity test
turned ``slope is positive'' into ``function is increasing'' using the MVT. But the
MVT itself rested on Rolle, which rested on the Extreme Value Theorem, which rested
on the \emph{completeness} of the real numbers --- the fact that there are no gaps in
the number line. So when you say ``a positive derivative makes a function climb,''
you are, four layers down, relying on there being no holes in $\R$. Can you feel why
a function on the rationals only, with no irrational points, could fail the whole
chain? (Hint: a continuous function on the rationals can fail to attain its
maximum, because the max might want to live at an irrational point that isn't
there.)
\end{walkbox}

\section{The Cauchy Mean Value Theorem}

There is a powerful generalization that you will need the moment you reach
l'H\^opital's rule, so we prove it here while the machinery is fresh. It compares the
changes in \emph{two} functions over the same interval.

\begin{theorem}[Cauchy's Mean Value Theorem]\label{ch15:thm-cauchy}
Let $f$ and $g$ be continuous on $[a,b]$ and differentiable on $(a,b)$. Then there
is a point $c\in(a,b)$ with
\[
  \bigl(f(b)-f(a)\bigr)\,g'(c)=\bigl(g(b)-g(a)\bigr)\,f'(c).
\]
If $g'$ is never zero on $(a,b)$ (so $g(b)\ne g(a)$), this rearranges to
\[
  \frac{f(b)-f(a)}{g(b)-g(a)}=\frac{f'(c)}{g'(c)}.
\]
\end{theorem}

\begin{intuitionbox}
Think of $\bigl(g(t),f(t)\bigr)$ as a point tracing a curve in the plane as $t$ runs
from $a$ to $b$. The ordinary MVT compared a slope to an average slope along the
$x$-axis. Cauchy's version compares the \emph{direction the curve is heading},
$f'(c)/g'(c)$, to the \emph{straight-line direction} from start to end,
$\frac{f(b)-f(a)}{g(b)-g(a)}$. It is the MVT for a parametrized path: somewhere the
path points the same way as the chord joining its endpoints.
\end{intuitionbox}

\begin{proof}
We use the same subtract-a-helper trick that turned Rolle into the MVT, only now the
helper is chosen to kill both endpoints at once. Define
\[
  h(x)=\bigl(f(b)-f(a)\bigr)\,g(x)-\bigl(g(b)-g(a)\bigr)\,f(x).
\]
This $h$ is continuous on $[a,b]$ and differentiable on $(a,b)$. Compute the
endpoint values:
\begin{align*}
  h(a)&=\bigl(f(b)-f(a)\bigr)g(a)-\bigl(g(b)-g(a)\bigr)f(a)
      =f(b)g(a)-g(b)f(a),\\
  h(b)&=\bigl(f(b)-f(a)\bigr)g(b)-\bigl(g(b)-g(a)\bigr)f(b)
      =f(b)g(a)-g(b)f(a).
\end{align*}
The cross terms cancel and $h(a)=h(b)$. So Rolle applies: there is $c\in(a,b)$ with
$h'(c)=0$. But
\[
  h'(x)=\bigl(f(b)-f(a)\bigr)g'(x)-\bigl(g(b)-g(a)\bigr)f'(x),
\]
and $h'(c)=0$ is precisely the claimed identity.
\end{proof}

\begin{connectionbox}
Notice the ordinary MVT is the special case $g(x)=x$: then $g'(c)=1$ and
$g(b)-g(a)=b-a$, and Cauchy's identity collapses to
$f(b)-f(a)=f'(c)(b-a)$. So Cauchy's theorem genuinely contains the Lagrange MVT, and
both descend from Rolle. Three theorems, one idea --- ``a difference that vanishes at
both ends has a flat spot inside'' --- applied to three cleverly chosen helper
functions. When you reach l'H\^opital's rule a couple of chapters from now, that
quotient $\frac{f'(c)}{g'(c)}$ will be exactly what lets a ratio of functions be
replaced by a ratio of their derivatives.
\end{connectionbox}

\section{A glance ahead: the MVT seeds Taylor's theorem}

The increment form $f(b)-f(a)=f'(c)(b-a)$ can be read as a \emph{first-order} exact
remainder: the change in $f$ is captured exactly by one derivative at one mystery
point. Taylor's theorem, which you will meet in full later, is the same idea pushed
to higher order. It says
\[
  f(b)=f(a)+f'(a)(b-a)+\tfrac12 f''(c)(b-a)^2
\]
for some interior $c$ --- a ``second-order MVT'' whose proof is, once again, build a
helper that vanishes at both ends and apply Rolle. The whole tower of polynomial
approximation, and every error bound that comes with it, grows from the seed you
planted in this chapter.

\begin{aestheticbox}
Step back and look at what one humble picture accomplished. ``An average slope is
attained somewhere'' gave us: the meaning of the constant of integration, the test
that turns derivative signs into shapes, a rigorous reason a zero rate keeps a
quantity fixed, sharp numerical error bounds, a clean way to prove inequalities, the
key lemma behind l'H\^opital, and the seed of Taylor's theorem. Few facts in
mathematics carry so much on so slender a frame. That is what we mean by a
\emph{workhorse}: it is rarely the star of the show, but almost nothing downstream
would stand without it.
\end{aestheticbox}

\section{Exercises}

\begin{exercise}[Finding $c$]
For each function and interval, verify the MVT hypotheses and find every $c$ the
theorem promises.
(a) $f(x)=x^2-3x$ on $[1,4]$.\quad
(b) $f(x)=x+\frac1x$ on $[1,3]$.\quad
(c) $f(x)=\sin x$ on $[0,\pi]$.
\end{exercise}

\begin{exercise}[Rolle by hand]
Let $f(x)=x^3-3x$. Find two values $a<b$ with $f(a)=f(b)$, then locate the interior
point(s) where $f'=0$ guaranteed by Rolle's theorem. Sketch the graph and mark the
horizontal tangents.
\end{exercise}

\begin{exercise}[Why a hypothesis is needed]
Give an explicit function on $[-1,1]$ that is continuous, has $f(-1)=f(1)$, but has
\emph{no} point where $f'=0$. Explain precisely which hypothesis of Rolle's theorem
your example violates, and where.
\end{exercise}

\begin{exercise}[A speed argument]
A car travels $90$ miles in exactly $90$ minutes. Use the MVT to prove that at some
instant its instantaneous speed was exactly $60$ mph. Then prove the sharper claim
that at some instant its speed was \emph{at least} $60$ mph, using only the average,
without the MVT. Which statement is stronger, and why does the MVT earn the
stronger one?
\end{exercise}

\begin{exercise}[Bounding with the derivative]
Suppose $f$ is differentiable on $\R$ with $f(1)=5$ and $\abs{f'(x)}\le 2$ for all
$x$. Find the smallest interval that must contain $f(4)$. (Use
Corollary~\ref{ch15:cor-bound}.)
\end{exercise}

\begin{exercise}[Inequalities from monotonicity]
Prove each inequality by forming a difference $g$, signing $g'$, and using the
monotonicity test.
(a) $\tan x > x$ for $0<x<\frac\pi2$.\quad
(b) $\cos x > 1-\frac{x^2}{2}$ for $x\ne 0$.\quad
(c) $\frac{x}{1+x}<\ln(1+x)$ for $x>0$.
\end{exercise}

\begin{exercise}[Uniqueness of a solution]
Show that the equation $x^5+x-1=0$ has exactly one real solution. (Existence: use
the Intermediate Value Theorem on $[0,1]$. Uniqueness: if there were two solutions,
Rolle's theorem would force $f'=0$ somewhere between them --- show that is
impossible.)
\end{exercise}

\begin{exercise}[Constant on an interval]
Suppose $f$ is differentiable on $\R$ with $f'(x)=3$ for all $x$. Prove, using the
MVT and not antidifferentiation rules, that $f(x)=3x+f(0)$ for every $x$.
\end{exercise}

\begin{exercise}[A fixed-point bound]
Let $g$ be differentiable on $\R$ with $\abs{g'(x)}\le \tfrac12$ for all $x$. Show
that for any two numbers $p,q$, $\abs{g(p)-g(q)}\le \tfrac12\abs{p-q}$. (Such a $g$
is a \emph{contraction}; this inequality is the heart of why iterating $g$ converges
to a single fixed point.)
\end{exercise}

\begin{exercise}[Proof: a derivative that is always positive]
Prove carefully, citing the MVT, that if $f'(x)>0$ for all $x$ in an open interval
$I$, then $f$ is one-to-one on $I$. Where exactly do you use that the interval is
connected?
\end{exercise}

\begin{exercise}[Cauchy MVT in action]
Let $f(x)=x^2$ and $g(x)=x^3$ on $[1,2]$. Find the point $c$ guaranteed by Cauchy's
Mean Value Theorem, i.e. solve $\frac{f(2)-f(1)}{g(2)-g(1)}=\frac{f'(c)}{g'(c)}$ for
$c\in(1,2)$. Confirm your $c$ lies in the interval.
\end{exercise}

\begin{exercise}[Conceptual: the role of the EVT]
Explain in a short paragraph why Rolle's theorem genuinely needs the Extreme Value
Theorem, and why the Extreme Value Theorem in turn needs the completeness of $\R$.
Give an example of a continuous function on the rationals in $[0,2]$ (with the usual
formula) that does \emph{not} attain a maximum, illustrating what breaks without
completeness.
\end{exercise}
